%% LabSetup

%% Automated Container Environment
\section{Automated Building of IBE Lab environment}
\label{apx:lab-final_envbuild}


Podman is a daemonless, open source, Linux native tool designed to make it easy to find,
run, build, share and deploy applications
using Open Containers Initiative (OCI)\cite{oci.containers} Containers and
Container Images.\cite{podman} Similarly to other common Container Engines (Docker, CRI-O, containerd),
Podman relies on an OCI compliant Container Runtime (\lstinline{runc}, \lstinline{crun}, \lstinline{runv}, etc.)
to interface with the operating system and create the running containers.\cite{podman}

\subsection{Cleanup the existing environment}

The \lstinline{IBE_LAB.yml} will build the whole lab environment from scratch,
thus any previously made network and containers related to IBE lab can be removed
by running the cleanup script showed within \textit{\ref{apx.code.env.cleaup} snippet}.

%% TODO: Snippet from `cleanup_ibe.lab.sh`
\begin{lstlisting}[language=bash, caption=Host preparation for building the IBE environment, label=apx.code.env.cleaup]
    # IBE specific cleanup script snippet <-
\end{lstlisting}

% ---

\subsection{Review of \textit{Dockerfiles}\cite{docker.dockerfiles}}
\label{apx:lab-final_dockerfiles}

A \lstinline{Dockerfile} is basically a text document that contains all the necessary
commands to call on the command line to assemble an image.


% TODO; A vegleges architecture alapjan
For this specific lab I assembled the files according the network topology:
\lstinline{Clientfile}, \lstinline{Serverfile},
\lstinline{SiteA}, \lstinline{SiteB}


\subsubsection{Serverfile}
\label{apx:dockerfile.serverfile}

\begin{lstlisting}[language=bash,caption=Serverile overview for container vpn-server-10.9.0.11]
# ===========================================================